% GENERATED by paper_cvpr/convert.py from paper/main.tex -- do not edit.

\section{Provenance of every column}
\label{app:provenance}
The audit is only as strong as its coverage, so we state it rather than summarising it. Each
column of Tab.~\ref{tab:battery} of the main paper rests on the following runs and cells, emitted alongside the
table by the same script.

\begin{table*}[t]
\centering
\small
\begin{tabular}{llrl}
\toprule
benchmark & run directory & cells & seeds $\times$ object counts \\
\midrule
motion & \texttt{onset\_shortcut\_audit\_5seed} & 15 & seeds 0,1,2,3,4 $\times$ N=3,5,8 \\
onset & \texttt{onset\_shortcut\_audit\_5seed} & 15 & seeds 0,1,2,3,4 $\times$ N=3,5,8 \\
interaction & \texttt{audit\_interaction\_billiards} & 6 & seeds 0,1,2 $\times$ N=5,8 \\
clutter-onset & \texttt{audit\_onset\_clutter12k} & 6 & seeds 0,1,2 $\times$ N=5,8 \\
\bottomrule
\end{tabular}

\end{table*}

\paragraph{Split sizes.} Rows per cell after filtering, which is the quantity that decides
whether a cell can be trained on at all: the motion and onset benchmarks sit near
$750$--$950$ training rows, and the clutter-onset benchmark at $12{,}000$ episodes per cell
reaches $697$--$1225$ training and $81$--$156$ test rows. We hold every cell to a floor of
$250$ training and $35$ test rows; no reported cell falls below it. Two constructions were
excluded on this ground rather than quietly included: $N=3$ for the interaction benchmark, and
$N=3$ for clutter-onset, which would need roughly $27{,}000$ episodes per cell to clear the
floor. Below that floor a learned model's score is as likely to reflect data starvation as
anything about the benchmark.

\section{The full battery}
\label{app:battery}
Tab.~\ref{tab:battery} of the main paper shows the rules the argument turns on. For completeness, here is every
rule the battery scores, on every benchmark, grouped as in Sec.~\ref{sec:battery}. The
fitted radii and the rule definitions are emitted alongside these numbers by the audit script.

Two things are visible here that the main table has no room for. First, the geometric group
degrades gracefully rather than collapsing: \texttt{pusher\_approaching} and
\texttt{two\_nearest\_to\_pusher} hold $0.66$--$0.73$ almost everywhere, so a benchmark that
defeats one geometric rule has usually not defeated the family. Second,
\texttt{second\_nearest\_to\_pusher} is the only rule that is bad on every benchmark
($0.10$--$0.41$), which is the sanity check on the battery: a proximity rule pointed at the
wrong object should fail, and it does.

\begin{table*}[t]
\centering
\small
\begin{tabular}{llrrrr}
\toprule
rule & group & motion & onset & interaction & clutter-onset \\
\midrule
\texttt{already\_moving} & kinematic & 0.866 & 0.182 & 0.783 & 0.449 \\
\midrule
\texttt{nearest\_to\_pusher} & geometric & 0.836 & 0.928 & 0.437 & 0.656 \\
\texttt{pusher\_approaching} & geometric & 0.529 & 0.691 & 0.729 & 0.683 \\
\texttt{pusher\_near} & geometric & 0.812 & 0.932 & 0.715 & 0.698 \\
\texttt{second\_nearest\_to\_pusher} & geometric & 0.163 & 0.095 & 0.405 & 0.219 \\
\texttt{two\_nearest\_to\_pusher} & geometric & 0.692 & 0.698 & 0.679 & 0.656 \\
\midrule
\texttt{moving\_or\_near\_mover} & propagated & 0.866 & 0.302 & 0.955 & 0.579 \\
\texttt{near\_a\_mover} & propagated & 0.263 & 0.210 & 0.712 & 0.441 \\
\texttt{near\_pusher\_or\_mover} & propagated & 0.831 & 0.935 & 0.820 & 0.782 \\
\midrule
\texttt{moving\_or\_near} & combination & 0.927 & 0.961 & 0.858 & 0.821 \\
\midrule
\texttt{always\_change} & floor & 0.413 & 0.387 & 0.673 & 0.485 \\
\midrule
\texttt{logistic\_on\_distance} & control & 0.810 & 0.919 & 0.731 & 0.577 \\
\midrule
learned gate (velocity features) & learned & 0.824 & 0.508 & 0.709 & 0.425 \\
learned gate (velocity-free) & learned & 0.688 & 0.703 & 0.772 & 0.338 \\
\bottomrule
\end{tabular}

\end{table*}

\section{Per-domain characterisation}
\label{app:domains}
The precondition for the velocity shortcut is measured directly on the simulators, with no model
trained: three object counts $\times$ three seeds $\times$ $60$ episodes per domain, produced by
\texttt{experiments/domain\_characterization.py}. The motion threshold for each domain is derived
as described in Sec.~\ref{sec:setup}.

\begin{table*}[t]
\centering
\small
\begin{tabular}{lllrrrrr}
\toprule
domain & engine & contact regime & $P(\text{chg}\mid\text{moving})$ & $P(\text{chg}\mid\text{rest})$
& ratio & changed frac. & threshold \\
\midrule
\texttt{billiards} & Box2D & near-elastic & $0.982$ & $0.020$ & $50\times$ & $0.394$ & $0.031$ \\
\texttt{planar} & ours & quasi-static & $0.950$ & $0.011$ & $84\times$ & $0.152$ & $0.010$ \\
\texttt{clutter} & Chipmunk2D & high-friction, packed & $0.863$ & $0.024$ & $36\times$ & $0.345$ & $0.028^{*}$ \\
\texttt{tabletop} & MuJoCo & impulsive 3D & $0.714$ & $0.015$ & $48\times$ & $0.124$ & $0.020$ \\
\bottomrule
\end{tabular}

\smallskip
\footnotesize $^{*}$The clutter benchmarks are built with a threshold of $0.029$.
\end{table*}

The shortcut exists in every domain. It is strongest under near-elastic contact (billiards,
$0.982$), where objects keep moving after a hit, and weakest on the MuJoCo tabletop ($0.714$),
whose 3D solver lets a box come to rest through tilting and settling that a 2D engine cannot
represent, so more of its moving objects stop within one step.

The threshold derivation has two self-checks. It returns $0.020$ for the tabletop, the value it
was anchored to, and $0.010$ for planar, the value that domain independently requires. The
anchor is the fraction of steps retained rather than a statistic of moved-object displacement,
since the latter is measured on an already-filtered split and would make the derivation
circular.

